\scanpage{57}
l'ens des \deleted{facettes} $y\in\operatorname{Fac}_{i\deleted{\unsure{-1|n}}}(C)$ qui majorent $x$ soit $d_i$. Donc dans ce cas on a
\[
\operatorname{card}\operatorname{Rep}(C)=d_{\deleted{\unsure{1}}0}d_{\unsure{1}}\ldots d_{\unsure{n|n-1}}=d_0\ldots d_{n-1}
\]
(car $d_n=1$). \deleted{On a alors, si $N=\operatorname{Card}\operatorname{Aut}(C)$}\deleted{$N_c=\operatorname{Card}\operatorname{Aut}(\operatorname{Fac}(C))$}

\unsure{Evidemt|Evidemment} régulier $\Rightarrow$ homogène

\deleted{$N\mid N_c\mid\operatorname{card}D$ Evidemment, \unsure{si}}

Si
\[
\begin{aligned}
\deleted{\unsure{R}}N&=\operatorname{Card}\operatorname{Aut}(C)\\
N_c&=\operatorname{Card}\operatorname{Aut}(\operatorname{Fac}(C))\\
R&=\operatorname{Card}\operatorname{Rep}(C),
\end{aligned}
\]
on aura
\[
N\mid N_c\mid R
\]
et $C$ régulier $\Longleftrightarrow N=R$ (donc $N=N_c$), donc les autom combinatoires sont euclidiens) Donc dans ce cas on a une décomposition
\[
N=N_c=d_0\ldots d_n
\]

\noindent\rule{2em}{0.4pt}\par\nobreak\noindent\rule{2em}{0.4pt}

Orientation d'un polyèdre \inserted{\unsure{conv.|convexe}} $C$ : par définition on \unsure{en|sera} une orientation de l'espace affine $E_C$ qu'il engendre. \inserted{L'ens. de ses o. est de ordre 2} Mais on montre que l'ens $\omega(C)$ ne dépend, à isom canonique près, que de $\Phi(C)=\operatorname{Fac}(C)$ (polyèdre combinatoire, strict associé) Donc il y a \deleted{Isom}
\[
\begin{aligned}
\operatorname{Isom}_{\mathrm{ord}}(\Phi(C),\Phi(C'))&\longrightarrow\operatorname{Isom}(\omega(C),\omega(C'))\\
\operatorname{Aut}_{\mathrm{ord}}(\Phi(C))&\longrightarrow\{\pm1\}\qquad\text{(hom de groupes)}\\
&\hspace{2em}\shortparallel\\[-3pt]
&\hspace{1em}\unsure{\operatorname{Aut}_{\mathrm{Ens}}|\operatorname{Aut}_{\mathrm{ens}}}(\omega(C))
\end{aligned}
\]
Cas où $C$ admet un antipodisme : sa sign. est $(-1)^n$
% archive sheet number: 56
